% ============================================================
%  DBMS_10 – Projektvorschlag für die Ausarbeitung
%  THGA Bochum · Datenbanken · Stephan Bökelmann
% ============================================================
\documentclass[11pt,a4paper]{article}

\usepackage{thga-db}

% ----- Metadaten -------------------------------------------
\renewcommand{\thgaDatum}{Sommersemester 2026}

\begin{document}
\thispagestyle{empty}
\thgaTitel{Übung 10 – Projektvorschlag}%
{Ausarbeitung: Ein eigenes datenbankgestütztes System entwerfen}

\begin{infobox}
\textbf{Modul:} Einführung in Datenbankmanagementsysteme · THGA Bochum \\
\textbf{Dozent:} Stephan Bökelmann · \href{mailto:sboekelmann@ep1.rub.de}{sboekelmann@ep1.rub.de} \\
\textbf{Repository:} \url{https://github.com/MaxClerkwell/DBMS_10} \\
\textbf{Voraussetzungen:} Vorlesungen 01--10, Übungen DBMS\_01--DBMS\_09 \\
\textbf{Art der Übung:} Konzeption und schriftliche Abgabe -- \emph{kein} Programmierauftrag \\
\textbf{Umfang:} 2--4 Seiten Projektvorschlag
\end{infobox}

\tableofcontents
\vspace{1em}

% ============================================================
\section{Worum es in dieser Übung geht}
% ============================================================

Diese Übung ist bewusst \textbf{anders aufgebaut} als die bisherigen. Sie
schreiben diesmal keinen Code und arbeiten keine Schritt-für-Schritt-Anleitung
ab. Stattdessen \textbf{entwerfen} Sie ein eigenes Projekt und reichen dazu einen
\textbf{Projektvorschlag} beim Dozenten ein.

Der Projektvorschlag ist die Grundlage Ihrer \textbf{Ausarbeitung} -- der
Prüfungsleistung, in der Sie über das Semester hinweg ein eigenes,
datenbankgestütztes System bauen. Bevor Sie mit dem Bau beginnen, klären wir
gemeinsam den \emph{Gegenstand}: Welche Anwendung wollen Sie umsetzen? Welche
Daten verwaltet sie? Welche Funktionen bietet sie an?

\begin{infobox}
\textbf{Ziel dieser Übung:} Reichen Sie einen 2--4-seitigen Projektvorschlag
ein, aus dem hervorgeht, \emph{welches} datenbankgestützte System Sie in der
Ausarbeitung bauen wollen und \emph{wie} es die Inhalte der Vorlesungen 01--10
aufgreift.
\end{infobox}

Der Vorschlag wird vom Dozenten \textbf{gegengelesen und freigegeben}. Erst
nach der Freigabe beginnen Sie mit der Umsetzung. So stellen wir vor dem
eigentlichen Aufwand sicher, dass Ihr Vorhaben den richtigen Umfang hat --
weder zu klein noch zu groß -- und alle Pflichtbausteine abdeckt.

% ============================================================
\section{Die Ausarbeitung im Überblick}
% ============================================================

In der Ausarbeitung bauen Sie ein \textbf{vollständiges, datenbankgestütztes
System für eine Anwendung Ihrer Wahl}. Das Thema ist frei: eine
Vereinsverwaltung, ein Rezept- und Einkaufsplaner, eine Bibliothek, ein
Ticketsystem, ein Lager für einen Modellbau-Verein, eine Trainings-App --
alles ist erlaubt, solange sich sinnvolle Daten und Beziehungen modellieren
lassen.

Ihre Anwendung soll die Kette abbilden, die wir über das Semester aufgebaut
haben: \emph{von der Datenmodellierung bis zur laufenden, bedienbaren
Anwendung.}

\subsection{Der verbindliche Technologie-Stack}

Das System muss aus den folgenden Schichten bestehen. Sie kennen jede Schicht
bereits aus den Vorlesungen und den vorangegangenen Übungen:

\begin{itemize}[leftmargin=1.4em]
  \item \textbf{PostgreSQL} als relationales Datenbankmanagementsystem
        -- Ihr sauber modelliertes Schema mit Schlüsseln und
        Integritätsbedingungen (VL 03--07).
  \item \textbf{FastAPI} als REST-Schicht in Python -- die Datenbank wird
        \emph{nie} direkt vom Frontend angesprochen, sondern immer über die API
        (VL 08).
  \item \textbf{Docker Compose} als Betriebsumgebung -- Datenbank und API
        laufen als Dienste, die mit \texttt{docker compose up} gemeinsam
        starten (VL 09).
  \item \textbf{Ein Frontend}, das die API benutzt. Ein \textbf{Python-Frontend}
        (z.\,B. tkinter wie in VL 10) ist die empfohlene Variante, weil Sie es
        bereits kennen. Sie dürfen aber auch ein anderes Frontend bauen
        (Weboberfläche, Kommandozeilenwerkzeug o.\,Ä.), solange es die API
        nutzt.
\end{itemize}

\begin{beispielbox}
\textbf{Die Architektur in einem Satz:}\\[4pt]
\centerline{%
  \textbf{Frontend} $\;\longrightarrow\;$ \textbf{FastAPI (REST)}
  $\;\longrightarrow\;$ \textbf{PostgreSQL}%
}\\[4pt]
\centerline{alles orchestriert durch \textbf{Docker Compose}}
\end{beispielbox}

\subsection{Was eine sinnvolle Größe ist}

Rechnen Sie mit einem Datenmodell von etwa \textbf{drei bis sechs Tabellen} mit
mindestens einer \texttt{N:M}-Beziehung, einer Handvoll API-Endpunkte (lesend
\emph{und} schreibend) und einem Frontend, das die wichtigsten dieser Endpunkte
bedient. Das ist genug, um alle gelernten Konzepte zu zeigen, und wenig genug,
um im Semester fertig zu werden.

% ============================================================
\section{Was aus den Vorlesungen einfließen soll}
% ============================================================

Ihr Vorschlag -- und später die Ausarbeitung -- soll erkennbar auf den Inhalten
der Vorlesungen \textbf{01--10} aufbauen. Die folgende Tabelle zeigt, welcher
Baustein aus welcher Vorlesung stammt und woran man ihn im Projekt erkennt.

\renewcommand{\arraystretch}{1.3}
\begin{tabularx}{\linewidth}{@{}p{1.5cm}p{4.6cm}X@{}}
\toprule
\textbf{VL} & \textbf{Thema} & \textbf{Woran man es im Projekt erkennt} \\
\midrule
01 & Was ist ein DBMS, Drei-Ebenen-Modell & Begründung, \emph{warum} ein
     relationales DBMS zum Vorhaben passt \\
02 & ER-Modellierung, Kardinalitäten & Eine ER-Skizze mit Entitäten,
     Beziehungen und Kardinalitäten \\
03 & Relationales Modell, Schlüssel & Primär- und Fremdschlüssel, referentielle
     Integrität \\
04 & Normalisierung, Normalformen & Ein Schema mindestens in 3NF, ohne
     Redundanzen \\
05 & SQL DDL/DML & \texttt{CREATE TABLE}-Skript mit Constraints; Seed-Daten \\
06 & SQL-Abfragen, JOINs, Aggregation & Abfragen hinter den API-Endpunkten
     (mind. ein JOIN, eine Aggregation) \\
07 & PostgreSQL im Betrieb & Ausführung auf PostgreSQL, Init-Skript, Datentypen \\
08 & Python, psycopg2, FastAPI & Die REST-Schicht mit lesenden und schreibenden
     Endpunkten \\
09 & Docker, Compose, .env & \texttt{docker-compose.yml} mit Datenbank + API,
     Secrets in \texttt{.env} \\
10 & Frontend-Grundlagen & Ein bedienbares Frontend, das die API nutzt \\
\bottomrule
\end{tabularx}

\begin{infobox}
\textbf{Was \emph{nicht} verpflichtend ist:} Die vertiefenden Inhalte aus den
Vorlesungen \textbf{10--12} sind \emph{keine} Pflichtbestandteile der
Ausarbeitung. Konkret sind \textbf{optional}:
\begin{itemize}[leftmargin=1.4em, topsep=2pt]
  \item das \emph{Verpacken} des Frontends als nativer Installer
        (\texttt{.deb}, Setup-\texttt{.exe}, \texttt{.dmg}, PyInstaller) -- VL 10/13,
  \item fortgeschrittenes Transaktions- und Isolationsverhalten (ACID,
        Isolation Levels, Sperren, MVCC) -- VL 11,
  \item Zeitreihendaten, InfluxDB und Grafana -- VL 12.
\end{itemize}
Ein \emph{funktionierendes} Frontend genügt; es muss nicht paketiert oder
installierbar sein. Wer diese Themen dennoch einbaut, darf das gerne als
freiwillige Erweiterung tun.
\end{infobox}

% ============================================================
\section{Der Projektvorschlag -- Ihre Abgabe}
% ============================================================

Reichen Sie ein PDF von \textbf{2--4 Seiten} ein, das die folgenden Abschnitte
enthält. Nutzen Sie die Überschriften als Vorlage.

\subsection{Anwendungsdomäne und Motivation}

Beschreiben Sie in wenigen Sätzen, \emph{welche} Anwendung Sie bauen und
\emph{welches Problem} sie löst. Wer benutzt sie, und wofür? Ein oder zwei
konkrete Nutzungsszenarien machen den Vorschlag greifbar.

\subsection{Datenmodell -- ER-Skizze}

Zeichnen Sie eine \textbf{ER-Skizze} (VL 02) mit den zentralen Entitätstypen,
ihren wichtigsten Attributen und den Beziehungen samt \textbf{Kardinalitäten}.
Eine saubere Handzeichnung als Foto ist völlig ausreichend. Achten Sie darauf,
dass mindestens eine \texttt{N:M}-Beziehung vorkommt.

\subsection{Vom Modell zum Schema}

Leiten Sie aus dem ER-Modell ein \textbf{Relationenschema} ab (VL 03). Nennen
Sie pro Tabelle die Spalten, den \textbf{Primärschlüssel} und die
\textbf{Fremdschlüssel}. Begründen Sie kurz, dass Ihr Schema frei von
Redundanzen ist (mindestens 3NF, VL 04) -- ein bis zwei Sätze genügen.

\subsection{API-Entwurf}

Skizzieren Sie die geplanten \textbf{FastAPI-Endpunkte} (VL 08). Eine Tabelle
reicht:

\renewcommand{\arraystretch}{1.2}
\begin{tabularx}{\linewidth}{@{}p{2cm}p{4cm}X@{}}
\toprule
\textbf{Methode} & \textbf{Pfad} & \textbf{Zweck} \\
\midrule
\texttt{GET}  & \texttt{/…}  & liest Daten (mit JOIN/Aggregation im Hintergrund) \\
\texttt{POST} & \texttt{/…}  & legt einen neuen Datensatz an \\
\multicolumn{3}{@{}l}{\itshape … mindestens ein lesender und ein schreibender Endpunkt} \\
\bottomrule
\end{tabularx}

Markieren Sie, hinter welchem Endpunkt eine \textbf{Abfrage mit JOIN} und wo
eine \textbf{Aggregation} steckt (VL 06).

\subsection{Frontend-Idee}

Beschreiben Sie kurz, wie das Frontend aussieht und welche Endpunkte es bedient
(VL 10). Ein grobes Wireframe oder eine Aufzählung der Bildschirme/Masken
genügt. Geben Sie an, ob Sie die empfohlene Python-Variante (tkinter) oder etwas
anderes nutzen -- und im zweiten Fall, was und warum.

\subsection{Betrieb mit Docker Compose}

Listen Sie die \textbf{Dienste} auf, die in Ihrer \texttt{docker-compose.yml}
laufen werden (mindestens \texttt{postgres} und \texttt{api}), und wo die
Zugangsdaten liegen (\texttt{.env}, VL 09).

\subsection{Umfang, Meilensteine und Risiken}

Schätzen Sie den Umfang ehrlich ein: Wie viele Tabellen und Endpunkte? Nennen
Sie \textbf{2--3 Meilensteine} (z.\,B. ``Schema + Seed-Daten stehen'', ``API
lesend/schreibend'', ``Frontend bedient die Kern-Endpunkte'') und \textbf{ein
bis zwei Risiken}, an denen es scheitern könnte.

% ============================================================
\section{Formale Anforderungen an die Abgabe}
% ============================================================

\begin{itemize}[leftmargin=1.4em]
  \item \textbf{Format:} ein PDF, 2--4 Seiten, mit den Abschnitten aus
        Abschnitt~4.
  \item \textbf{Abgabe:} per E-Mail an
        \href{mailto:sboekelmann@ep1.rub.de}{sboekelmann@ep1.rub.de}. Legen Sie
        das PDF zusätzlich in Ihr Projekt-Repository (etwa als
        \texttt{docs/proposal.pdf}), sobald Sie eines angelegt haben.
  \item \textbf{Betreff / Dateiname:} \texttt{Ausarbeitung-Vorschlag
        <Nachname>}.
  \item \textbf{Sprache:} Deutsch oder Englisch.
\end{itemize}

\begin{warnbox}
\textbf{Beginnen Sie noch nicht mit der Implementierung}, bevor der Vorschlag
freigegeben ist. Warten Sie die Rückmeldung des Dozenten ab -- sie kann Umfang
oder Datenmodell noch verschieben und erspart Ihnen so unnötige Arbeit.
\end{warnbox}

% ============================================================
\section{Annahmekriterien -- Checkliste}
% ============================================================

Ihr Vorschlag wird angenommen, wenn er die folgenden Punkte erfüllt. Nutzen Sie
die Liste als Selbstkontrolle vor der Abgabe.

\begin{itemize}[leftmargin=1.8em, label=$\square$]
  \item Anwendungsdomäne und Nutzen sind klar beschrieben.
  \item Eine ER-Skizze mit Kardinalitäten liegt vor; mindestens eine
        \texttt{N:M}-Beziehung ist enthalten.
  \item Ein Relationenschema mit Primär- und Fremdschlüsseln ist abgeleitet
        (3 bis 6 Tabellen, mind. 3NF).
  \item Mindestens ein lesender und ein schreibender API-Endpunkt sind
        vorgesehen; JOIN und Aggregation sind verortet.
  \item Das Frontend und die von ihm genutzten Endpunkte sind skizziert.
  \item Der Docker-Compose-Betrieb (postgres + api, \texttt{.env}) ist benannt.
  \item Umfang, Meilensteine und Risiken sind realistisch dargestellt.
  \item Der verbindliche Stack (PostgreSQL, FastAPI, Docker Compose, Frontend)
        ist vollständig abgedeckt.
\end{itemize}

% ============================================================
\section{Ein durchgängiges Kurzbeispiel}
% ============================================================

Damit Sie das erwartete Niveau einschätzen können, hier ein stark verkürztes
Beispiel für die Domäne \emph{``Vereins-Bibliothek''}. Ihr eigener Vorschlag ist
ausführlicher, dies zeigt nur die Richtung.

\begin{beispielbox}
\textbf{Domäne.} Eine kleine Vereins-Bibliothek verwaltet Bücher, Mitglieder
und Ausleihen. Mitglieder sollen sehen, welche Bücher verfügbar sind, und der
Bibliothekar soll Ausleihen ein- und zurückbuchen.

\medskip
\textbf{Datenmodell (Ausschnitt).}
\begin{itemize}[leftmargin=1.4em, topsep=2pt]
  \item \texttt{mitglied}(\underline{id}, name, email)
  \item \texttt{buch}(\underline{id}, titel, autor, isbn, bestand)
  \item \texttt{ausleihe}(\underline{id}, \emph{mitglied\_id}, \emph{buch\_id},
        von, bis, zurueck) -- realisiert die \texttt{N:M}-Beziehung zwischen
        Mitglied und Buch.
\end{itemize}

\textbf{API (Ausschnitt).}
\begin{itemize}[leftmargin=1.4em, topsep=2pt]
  \item \texttt{GET /buecher} -- alle Bücher mit Anzahl aktueller Ausleihen
        (\emph{JOIN + Aggregation}).
  \item \texttt{GET /mitglieder/\{id\}/ausleihen} -- offene Ausleihen eines
        Mitglieds (\emph{JOIN}).
  \item \texttt{POST /ausleihen} -- Buch ausleihen.
  \item \texttt{POST /rueckgaben} -- Buch zurückbuchen.
\end{itemize}

\textbf{Frontend.} tkinter-Fenster mit zwei Reitern: ``Bücher'' (Tabelle +
Suchfeld) und ``Ausleihe'' (Formular für Ausleihe/Rückgabe).

\medskip
\textbf{Betrieb.} \texttt{docker-compose.yml} mit den Diensten \texttt{postgres}
und \texttt{api}; Passwörter in \texttt{.env}; Schema und Seed-Daten über ein
Init-Skript.

\medskip
\textbf{Meilensteine.} (1)~Schema + Seed-Daten; (2)~API lesend/schreibend;
(3)~Frontend bedient beide Reiter. \textbf{Risiko:} korrektes Zurückbuchen bei
gleichzeitig mehreren offenen Ausleihen desselben Buches.
\end{beispielbox}

% ============================================================
\section{Nächste Schritte}
% ============================================================

\begin{enumerate}[leftmargin=1.6em]
  \item Wählen Sie eine Domäne, die Sie interessiert -- Sie werden ein ganzes
        Semester daran arbeiten.
  \item Skizzieren Sie das ER-Modell auf Papier und leiten Sie das Schema ab.
  \item Füllen Sie die Abschnitte aus Abschnitt~4 aus und prüfen Sie die
        Checkliste in Abschnitt~6.
  \item Senden Sie das PDF an den Dozenten und warten Sie die Freigabe ab.
\end{enumerate}

\vspace{1em}
\begin{infobox}
\textbf{Weiterführend:} Die Übungen DBMS\_07 (PostgreSQL), DBMS\_08 (FastAPI),
DBMS\_09 (Docker Compose) und die Vorlesung~10 (Frontend) sind Ihre wichtigsten
Referenzen für die Umsetzung. Sie müssen für den Vorschlag nichts Neues lernen
-- nur das Bekannte auf Ihre eigene Idee anwenden.
\end{infobox}

\end{document}
